\documentclass[11pt]{article}
\usepackage[a4paper,margin=28mm]{geometry}
\usepackage[T1]{fontenc}
\usepackage{lmodern}
\usepackage{amsmath,amssymb,amsthm}
\usepackage{microtype}
\usepackage[colorlinks=true,linkcolor=blue,citecolor=blue,urlcolor=blue]{hyperref}
\newcommand{\T}{\mathbb T}
\numberwithin{equation}{section}
\newtheorem{theorem}{Theorem}[section]
\newtheorem{proposition}[theorem]{Proposition}
\newtheorem{lemma}[theorem]{Lemma}
\title{Uniform vanishing-viscosity convergence\\
for a local quadratic mean field game}
\author{Kaiwen Jin\\[2mm]
\small\href{mailto:yui.ui.siki@gmail.com}{\texttt{yui.ui.siki@gmail.com}}}
\date{15 October 2026}
\begin{document}
\maketitle

\begin{abstract}
We prove uniform vanishing-viscosity convergence for the local mean field
game system with quadratic Hamiltonian and linear density coupling on a
flat torus in dimensions one, two and three. For every smooth nonnegative
probability initial density, smooth fixed terminal cost and positive
time horizon, the entire family of value functions converges uniformly
on each interior time slab to a single continuous function.
Initial vacuum is allowed. Averaging finite-energy perturbations of
optimal trajectories turns a density-gradient estimate into a modulus
of continuity independent of viscosity. A convex variational limit and
path calibration then identify the occupied-region value, while
homogeneous Hamilton--Jacobi comparison determines its vacuum extension.
\end{abstract}

\section{Introduction}\label{sec:intro}

Mean field games with local coupling describe a population whose individual
running cost depends on the density at the current position. In the
quadratic case, a backward Hamilton--Jacobi equation for the value function
is coupled to a forward Fokker--Planck equation for the population density.
Positive viscosity represents diffusion in the individual trajectories.
For smooth periodic data, classical solutions at each positive viscosity
are available \cite[Theorem~4.1]{CLLP}. The question here concerns their
deterministic limit, rather than their existence at fixed viscosity.

Weak formulations of first-order mean field games with local coupling
can be obtained from convex variational problems \cite{Cardaliaguet}.
Spatial Sobolev estimates for such variational systems are developed in
\cite{GraberMeszaros}. These methods control the density and value
gradient, but uniform convergence of values also requires control of
their behavior where the density is small. Tang and Zhang
\cite[Section~7, problem~2(a)]{TangZhang} ask whether the value
functions converge locally uniformly as viscosity vanishes. Their
local-coupling convergence estimates include density-weighted value
bounds, which do not alone control values in vacuum.

We treat the quadratic Hamiltonian with linear local coupling and a
fixed terminal cost. Let \(\T^d=\mathbb R^d/\mathbb Z^d\), with volume
one. For smooth data \(m_0,g\) and \(\nu>0\), consider
\begin{equation}\label{eq:mfg}
\begin{aligned}
 -\partial_tu_\nu-\nu\Delta u_\nu+\tfrac12|Du_\nu|^2&=m_\nu,\\
 \partial_tm_\nu-\nu\Delta m_\nu-\operatorname{div}(m_\nu Du_\nu)&=0
                &&\text{in }(0,T)\times\T^d,
\end{aligned}
\end{equation}
with
\begin{equation}\label{eq:data}
 m_\nu(0)=m_0,\qquad u_\nu(T)=g .
\end{equation}
The initial density may vanish. Uniform convergence is therefore a
statement about values on the whole torus, including regions that
density-weighted estimates may not see.

\begin{theorem}\label{thm:main}
Let \(d\in\{1,2,3\}\), \(T>0\),
\(m_0\in C^\infty(\T^d;[0,\infty))\) with \(\int_{\T^d}m_0=1\),
and \(g\in C^\infty(\T^d;\mathbb R)\).
Let \((u_\nu,m_\nu)\) be the classical solutions of
\eqref{eq:mfg}--\eqref{eq:data}. There is a function
\(u\in C((0,T)\times\T^d)\) such that, for every \(0<\tau<T/2\),
\begin{equation}\label{eq:main-limit}
 \|u_\nu-u\|_{L^\infty([\tau,T-\tau]\times\T^d)}
       \longrightarrow0\qquad\text{as }\nu\downarrow0 .
\end{equation}
The limit is the same for the entire family of viscosities.
\end{theorem}

Thus the uniform-convergence question has an affirmative answer for this
quadratic, linear-coupling specialization in dimensions at most three,
without a positive lower bound on the initial density or a restriction on
the time horizon. The theorem does not address general Hamiltonians,
other couplings or higher dimensions. We do not obtain a convergence rate.

The first difficulty is to obtain continuity estimates without a uniform
pointwise bound on the density. A differentiated identity bounds its
spatial gradient in terms of the value bound. Averaging finite-energy
paths that spread from a point closes this estimate. To compare two
nearby starting positions, we perturb an optimal trajectory by an
independent family of shifts which return to the original trajectory
after a short time. The averaged running cost is controlled by the
density's space-time norm, and the kinetic increment gives a modulus
of continuity independent of viscosity.

The second difficulty is that occupied-region information does not
immediately determine a value in vacuum. A heat-semigroup recovery
sequence preserves the initial density exactly and identifies a common
first-order density and flux. Their finite-action path law calibrates
every continuous value limit on the occupied region. On an open region
where two limits differ, the common density must be zero; there
homogeneous Hamilton--Jacobi comparison proves that the limits coincide.
This uses the variational and calibration ideas of
\cite{Cardaliaguet}, with the required specialization proved below.

Section~\ref{sec:estimates} gives the uniform estimates and compactness.
Section~\ref{sec:limit} identifies the density limit and proves
calibration. Section~\ref{sec:vacuum} completes the value comparison.
All spatial integrations use the normalized torus measure; \(D\) and
\(\Delta\) act only in space, and \(\operatorname{dist}\) denotes the
shortest-lift torus distance.
\section{Uniform estimates}\label{sec:estimates}


For clarity, the fixed-viscosity classical theory cited above applies in
the present normalization. Set \(s=\nu t\),
\(U(s,x)=u_\nu(s/\nu,x)/\nu\) and \(n(s,x)=m_\nu(s/\nu,x)\).
The equations have unit diffusivity, Hamiltonian \(|DU|^2/2\),
coupling \(n/\nu^2\), horizon \(\nu T\), terminal value \(g/\nu\),
and initial density \(m_0\). The coupling is smooth and bounded below
on nonnegative densities, so \cite[Theorem~4.1 and its regularity
conclusion]{CLLP} applies. Smooth data permit the differentiations and
endpoint integrations used below. This is only qualitative regularity
for each fixed \(\nu\).

Throughout the proof \(0<\nu\le1\), \(Q=(0,T)\times\T^d\), and \(D\)
denotes the spatial gradient. Constants may depend on \(d,T,m_0,g\), and
constants denoted \(C_\tau\) may also depend on \(\tau\). They never depend
on \(\nu\). Mass conservation and positivity give
\[
 m_\nu\ge0,\qquad \int_{\T^d}m_\nu(t,x)\,dx=1 .
\]
These facts follow by integrating the forward equation and testing its
negative part. The resulting inequality
\(\frac d{dt}\|m_\nu^-\|_2^2\le C_\nu\|m_\nu^-\|_2^2\), with zero
initial value, gives the latter assertion.

\subsection{Control verification and the density estimate}

We first record the control formula used to compare values at nearby points.
Let \(B\) be a standard \(d\)-dimensional Brownian motion. An admissible
control is an adapted measurable process \(a\) with
\(\mathbb E\int_t^T|a_s|^2\,ds<\infty\). Positions are projected onto the
torus from
\[
 X_s=x+\int_t^s a_r\,dr+\sqrt{2\nu}(B_s-B_t).
\]
It\^o's formula, with the periodic lift of \(u_\nu\), gives
\begin{equation}\label{eq:verification}
 \mathbb E\left[g(X_T)+\int_t^T
       \left(\tfrac12|a_s|^2+m_\nu(s,X_s)\right)ds\right]
 =u_\nu(t,x)+\tfrac12\mathbb E\int_t^T
       |a_s+Du_\nu(s,X_s)|^2ds .
\end{equation}
The martingale has zero expectation because the test gradient is bounded
for fixed \(\nu\); all drift terms are integrable by Cauchy--Schwarz.
The coefficient \(\sqrt{2\nu}\) gives precisely the term \(\nu\Delta u_\nu\).
See, for example, \cite[Theorem~3.7.1]{Lawler} for this version of
It\^o's formula.

The smooth feedback \(a_s=-Du_\nu(s,X_s)\) is admissible and realizes
equality. Its process can be constructed by Picard iteration of the
integral equation for \(X-\sqrt{2\nu}B\), successively on short time
intervals. Thus \(u_\nu\) is the infimum of the left side of
\eqref{eq:verification}. Applying the same calculation on \([t,s]\),
with terminal function \(u_\nu(s,\cdot)\), gives
\begin{equation}\label{eq:interval}
 u_\nu(t,x)=\inf_a\mathbb E\left[
 u_\nu(s,X_s)+\int_t^s
 \left(\tfrac12|a_r|^2+m_\nu(r,X_r)\right)dr\right].
\end{equation}
An independent auxiliary random variable merely averages admissible costs
and therefore also supplies an upper bound for these infima.

Set \(g_*=\min g\) and
\(M_\nu=\max_{[0,T]\times\T^d}(u_\nu-g_*)\).
Nonnegative running costs imply \(u_\nu\ge g_*\). Every optimal trajectory
starting at any \((t,x)\) satisfies
\begin{equation}\label{eq:optimal-energy}
 \mathbb E\int_t^T|a_s|^2ds\le2M_\nu .
\end{equation}

The density gradient can be bounded in terms of this value bound.
For \(u_i=\partial_i u_\nu\), \(m_i=\partial_i m_\nu\), differentiating
the equations gives
\[
 u_{i,t}=-\nu\Delta u_i+Du_\nu\cdot Du_i-m_i,\qquad
 m_{i,t}=\nu\Delta m_i+
 \operatorname{div}(m_iDu_\nu+m_\nu Du_i).
\]
Consequently, periodic integration by parts yields
\[
 \frac d{dt}\int_{\T^d}m_i u_i
   =-\int_{\T^d}m_i^2-\int_{\T^d}m_\nu|Du_i|^2 .
\]
The diffusion terms cancel, as do the two terms containing
\(m_iDu_\nu\cdot Du_i\). Sum in \(i\) and integrate in time:
\begin{equation}\label{eq:derivative-energy}
 \int_Q\bigl(|Dm_\nu|^2+m_\nu|D^2u_\nu|^2\bigr)
 =-\int_{\T^d}\Delta m_0\,u_\nu(0)
       +\int_{\T^d}m_\nu(T)\Delta g
 \le \|\Delta m_0\|_1 M_\nu+\|\Delta g\|_\infty .
\end{equation}
Here \(\int\Delta m_0=0\) removes the constant \(g_*\). No division by
the density is involved.

For later use, the mass-one Gagliardo--Nirenberg estimate on the torus is
\begin{equation}\label{eq:gn}
 \int_{\T^d}m^{p_d}\le C_d(1+\|Dm\|_2^2),
 \qquad m\ge0,\quad\int_{\T^d}m=1,\quad
 p_d=\frac{2(d+1)}d .
\end{equation}
To see the additive term explicitly, multiply the periodic extension by
a fixed smooth compact cutoff and apply the Euclidean inequality
of \cite[Lecture~II, (2.2)--(2.3), pp.~125--126]{Nirenberg}.
Writing \(Y=\|m\|_{p_d}\), \(D_0=\|Dm\|_2\),
\(\theta=d/(d+1)\), one obtains
\(Y\le C(D_0+\|m\|_2)^\theta\).
Interpolation with mass one gives
\(\|m\|_2\le Y^{(d+1)/(d+2)}\).
Absorbing \(CY^{d/(d+2)}\) therefore gives
\(Y\le C(1+D_0^\theta)\); since \(\theta p_d=2\), this proves
\eqref{eq:gn}. The exceptional Euclidean condition in dimension two
requires \(\theta<1\), which holds. The additive constant is necessary
even for \(m=1\).

Since \(p_d\ge8/3\) in the dimensions under consideration,
\eqref{eq:derivative-energy} and \eqref{eq:gn} give, with \(p=8/3\),
\begin{equation}\label{eq:mp-before}
 \|m_\nu\|_{L^p(Q)}\le C(1+M_\nu)^{1/p}.
\end{equation}
We close this estimate by averaging paths that spread from their initial
point. Choose
\[
 \beta=\frac9{17},\qquad c_T=\min(1,T^{-\beta}).
\]
The conditions needed are \(\beta>1/2\) for finite kinetic energy and
\(d\beta<p-1\) for integrability of the running cost. They both hold
for \(d\le3\). If \(h=T-t\), take \(z\) uniform in the Euclidean unit
ball, independent of \(B\), and use
\[
 a_{t+r}=c_T\beta r^{\beta-1}z,\qquad
 X_{t+r}=x+c_Tr^\beta z+\sqrt{2\nu}B_r,\quad 0<r<h.
\]
Its kinetic cost is at most \(C_T h^{2\beta-1}\).
For \(F\ge0\), the ball average satisfies
\begin{equation}\label{eq:ball}
 \mathbb E_z F(Y+\vartheta z)
   \le C_d\vartheta^{-d/p}\|F\|_{L^p(\T^d)},\qquad
           0<\vartheta\le1 .
\end{equation}
Indeed, its periodic density is bounded by \(C_d\vartheta^{-d}\),
with support volume at most \(C_d\vartheta^d\); the overlap of periodic
lifts is bounded. H\"older's inequality gives \eqref{eq:ball}.
Conditioning on the Brownian position and then applying H\"older in time,
\begin{align}
 \mathbb E\int_0^h m_\nu(t+r,X_{t+r})dr
 &\le C_T
   \left(\int_0^h r^{-d\beta/(p-1)}dr\right)^{1/p'}
       \|m_\nu\|_{L^p(Q)} \notag\\
 &\le C_T h^{a_d}\|m_\nu\|_{L^p(Q)},\qquad
       a_d=1-\frac{d\beta+1}{p}>0 .\label{eq:spread-cost}
\end{align}
The worst value is \(a_3=1/34\).
The control formula, bounded terminal cost and \eqref{eq:mp-before}
now imply
\[
 M_\nu\le A+B(1+M_\nu)^{3/8}
\]
for fixed \(A,B\). Since \(3/8<1\), this bounds \(M_\nu\) uniformly.
We have proved
\begin{equation}\label{eq:uniform-bounds}
 \|u_\nu\|_\infty+\|m_\nu\|_{L^{8/3}(Q)}
 +\|Dm_\nu\|_{L^2(Q)}
 +\|\sqrt{m_\nu}D^2u_\nu\|_{L^2(Q)}\le C .
\end{equation}
Integrating the Hamilton--Jacobi equation also gives
\begin{equation}\label{eq:Du}
 \tfrac12\int_Q|Du_\nu|^2
      =T+\int_{\T^d}g-\int_{\T^d}u_\nu(0)\le C .
\end{equation}

\subsection{Spatial and temporal continuity}

The averaging argument can be arranged to return to an optimal trajectory
after a short time. This permits comparison of nearby initial positions
without changing the later cost.

\begin{proposition}\label{prop:modulus}
Let \(\alpha=1/34\). For every \(0<\tau<T\),
\begin{equation}\label{eq:modulus}
 |u_\nu(t,x)-u_\nu(s,y)|
 \le C_\tau\bigl(\operatorname{dist}(x,y)^\alpha
                         +|t-s|^{\alpha/2}\bigr),
 \quad t,s\in[0,T-\tau],\quad x,y\in\T^d .
\end{equation}
Moreover \(u_\nu(t,\cdot)\to g\) uniformly as \(t\uparrow T\),
uniformly in \(0<\nu\le1\).
\end{proposition}

\begin{proof}
First fix \(t\le T-\tau\). Put \(r=\operatorname{dist}(x,y)\) and
let \(q\) be a shortest lift of \(y-x\).
For \(0<r<\min(\tau/2,1/2)\), set \(h=r\), and let
\(\widehat X_\ell,\widehat a_\ell\) be an optimal process from \((t,x)\).
Write \(X_s=\widehat X_{t+s}\) and \(a_s=\widehat a_{t+s}\) for
elapsed time \(0\le s\le T-t\). Recall \(\beta=9/17\).
For an independent unit-ball parameter \(z\), define on \([0,h]\)
\[
 \vartheta_h(s)=h^\beta\left[\frac{s}{h}\left(1-\frac{s}{h}\right)\right]^\beta,
 \qquad
 X_s^z=X_s+(1-s/h)q+\vartheta_h(s)z .
\]
The new trajectory starts at \(y\), uses the same Brownian noise, and
equals \(X_h\) at time \(h\). Its adapted control is
\[
 a_s^z=a_s+b_s^z,\qquad b_s^z=-q/h+\vartheta'_h(s)z.
\]
After \(h\), use the original process and control. This is an admissible
control for each fixed \(z\): the added term is deterministic and
square-integrable. In fact,
\begin{equation}\label{eq:tube-energy}
 \int_0^h|b_s^z|^2ds
       \le C\left(\frac{r^2}{h}+h^{2\beta-1}\right).
\end{equation}
The derivative of \([v(1-v)]^\beta\) has squared endpoint power
\(2\beta-2=-16/17>-1\), which proves this estimate.

Since \(X_s\) is independent of \(z\), \eqref{eq:ball} applies to its
shifted position. The time integral now has two endpoints:
\[
 \int_0^h\vartheta_h(s)^{-d/(p-1)}ds
 =h^{1-d\beta/(p-1)}
   \int_0^1[v(1-v)]^{-d\beta/(p-1)}dv .
\]
It is finite because \(d\beta/(p-1)\le81/85<1\). With
\(p=8/3\) and \(a_d=1-(d\beta+1)/p\ge1/34\), this gives
\[
 \mathbb E\int_0^h m_\nu(t+s,X_s^z)ds
           \le C h^{a_d}\|m_\nu\|_{L^p(Q)} .
\]
All costs after \(h\) agree exactly. Subtract the original optimal
cost from the averaged competitor cost. Discard the original
nonnegative running cost, and bound the kinetic increment
\(a_s\cdot b_s^z+|b_s^z|^2/2\) by Cauchy--Schwarz,
\eqref{eq:optimal-energy}, \eqref{eq:uniform-bounds} and
\eqref{eq:tube-energy}. The result is
\[
 u_\nu(t,y)-u_\nu(t,x)
 \le C\left(\frac r{\sqrt h}+h^{\beta-1/2}
             +\frac{r^2}{h}+h^{2\beta-1}+h^{a_d}\right).
\]
With \(h=r\), every exponent is at least \(1/34\).
Swap \(x,y\), and use the uniform value bound for larger \(r\), to obtain
\begin{equation}\label{eq:space-modulus}
 |u_\nu(t,x)-u_\nu(t,y)|
       \le C_\tau\operatorname{dist}(x,y)^\alpha .
\end{equation}
The cost estimate concerns an averaged family of trajectories, rather
than a trace of an \(L^p\) function on a single curve.

For time continuity let \(0\le t<s\le T-\tau\), \(h=s-t\le1\).
In \eqref{eq:interval}, use the spreading control from
\eqref{eq:spread-cost}, restricted to this time interval.
Its kinetic and running costs are bounded by
\(C(h^{2\beta-1}+h^{a_d})\). The endpoint increment in
\eqref{eq:space-modulus} is bounded by
\[
 C_\tau\mathbb E\operatorname{dist}(X_s,x)^\alpha
 \le C_\tau(h^{\alpha\beta}+h^{\alpha/2})
 \le C_\tau h^{\alpha/2}.
\]
Here \(\nu\le1\) and the Brownian second moment were used.
This bounds \(u_\nu(t,x)-u_\nu(s,x)\) from above.
For the lower bound discard the nonnegative running cost in
\eqref{eq:interval}. If \(A=\mathbb E\int_t^s|a_r|^2dr\),
\[
 \mathbb E\operatorname{dist}(X_s,x)^\alpha
       \le C(hA)^{\alpha/2}+C(\nu h)^{\alpha/2}.
\]
Minimizing \(A/2-C_\tau h^{\alpha/2}A^{\alpha/2}\) over \(A\ge0\)
gives a lower bound \(-C_\tau h^{\alpha/(2-\alpha)}\).
Since \(\alpha/(2-\alpha)>\alpha/2\), the lower time increment
has the same modulus. Large time increments follow from the value
bound. Together with \eqref{eq:space-modulus}, this proves
\eqref{eq:modulus}.

Finally put \(h=T-t\le1\). The spreading control and the Lipschitz
continuity of \(g\) give
\begin{equation}\label{eq:terminal}
 g(x)-C(h+h^{1/2})\le u_\nu(t,x)
       \le g(x)+C(h^{2\beta-1}+h^{a_d}+h^{1/2}).
\end{equation}
The lower bound follows from the same energy minimization with the
Lipschitz exponent one. This proves the terminal assertion.
\end{proof}

Arzel\`a--Ascoli and a diagonal extraction on \([0,T-1/n]\times\T^d\)
now yield a locally uniform value limit for every vanishing sequence
of viscosities. In fact, \eqref{eq:terminal} makes the terminal strip
uniformly close to \(g\), so the extracted subsequence converges
uniformly on the whole closed cylinder. We will use its time-zero
trace when identifying the limit.

\section{The density limit and calibration}\label{sec:limit}

Compactness alone does not determine the value in vacuum. We first identify
the common density and flux by a convex minimization problem. A path
representation of that pair will then determine every value limit on the
occupied region.

\subsection{The convex action}

For \(\lambda\ge0\), let \(\mathcal K_\lambda\) consist of
narrowly continuous probability curves
\(\rho:[0,T]\to\mathcal P(\T^d)\), with a space-time density
\(\rho\in L^2(Q)\), and measurable fluxes \(k\) of finite kinetic action,
satisfying
\[
 \rho_t-\lambda\Delta\rho+\operatorname{div}k=0,\qquad \rho(0)=m_0 .
\]
Narrow continuity means continuity of the integral against each continuous
spatial test. The equation includes endpoint traces: for every smooth
\(\psi\) on the closed cylinder,
\begin{equation}\label{eq:weak-ce}
 \langle\psi(T),\rho(T)\rangle-\int_{\T^d}\psi(0)m_0
 =\int_Q \rho(\psi_t+\lambda\Delta\psi)+k\cdot D\psi .
\end{equation}
The density is specified at almost every time; the narrow curve supplies
the remaining time marginals. Define
\begin{equation}\label{eq:action}
 J_\lambda(\rho,k)=
 \int_Q\left(\frac{|k|^2}{2\rho}+\frac{\rho^2}{2}\right)
                    +\int_{\T^d}g\,d\rho(T),\qquad
 I_\lambda=\inf_{\mathcal K_\lambda}J_\lambda .
\end{equation}
The kinetic integrand is zero at \((\rho,k)=(0,0)\) and infinite at
\(\rho=0,k\ne0\). In particular, a finite-action flux vanishes in vacuum.
The terminal pairing allows a measure, rather than requiring a density at
every time.

\begin{lemma}\label{lem:primal}
Set \(j_\nu=-m_\nu Du_\nu\). Then \((m_\nu,j_\nu)\) minimizes
\(J_\nu\) and
\begin{equation}\label{eq:dual-value}
 I_\nu=\int_{\T^d}m_0u_\nu(0)-\frac12\int_Qm_\nu^2.
\end{equation}
Its density-square and kinetic integrals are bounded independently of
\(\nu\), and \(j_\nu\) is bounded in \(L^{4/3}(Q)\).
\end{lemma}

\begin{proof}
Test \eqref{eq:weak-ce} with \(u_\nu\), substitute
\(u_{\nu,t}+\nu\Delta u_\nu=|Du_\nu|^2/2-m_\nu\), and complete
two squares. For every finite-action competitor,
\begin{align}
 J_\nu(\rho,k)
 &=\int_{\T^d}m_0u_\nu(0)-\frac12\int_Qm_\nu^2
       +\frac12\int_Q(\rho-m_\nu)^2 \notag\\
 &\hspace{12mm}
       +\frac12\int_Q\frac{|k+\rho Du_\nu|^2}{\rho}.
       \label{eq:squares}
\end{align}
At vacuum the last term uses the same perspective convention.
The smooth test has bounded derivatives for fixed viscosity, so the weak
identity is legitimate for these competitors. Equation \eqref{eq:squares}
proves minimization and \eqref{eq:dual-value}.

Let \(P_s=e^{s\Delta}\) be the torus heat semigroup, \(P_0=\mathrm{id}\).
Its periodic Gaussian kernel is positive, symmetric and of mass one.
It commutes with spatial derivatives, preserves mass, and approximates
the identity uniformly on continuous functions. These properties follow
by summing the Euclidean Gaussian over \(\mathbb Z^d\); its derivative
series converge uniformly on compact positive-time intervals.
The competitor \((P_{\nu t}m_0,0)\) bounds
\[
 I_\nu\le \frac T2\|m_0\|_2^2+\max g .
\]
The terminal term in any admissible action is at least \(\min g\),
so the two nonnegative integral terms are bounded separately.
H\"older's inequality gives
\begin{equation}\label{eq:flux}
 \int_Q|j_\nu|^{4/3}
 \le\left(\int_Q\frac{|j_\nu|^2}{m_\nu}\right)^{2/3}
       \left(\int_Qm_\nu^2\right)^{1/3},
\end{equation}
which proves the flux bound.
\end{proof}

\begin{proposition}\label{prop:density}
There is a unique minimizer \((m,j)\) of \(J_0\) over \(\mathcal K_0\).
As \(\nu\downarrow0\),
\[
 I_\nu\longrightarrow I_0,\qquad
 m_\nu\longrightarrow m\quad\hbox{strongly in }L^2(Q).
\]
Along extracted subsequences the probability curves converge narrowly,
uniformly in time, to \(m(t)\), and \(j_\nu\rightharpoonup j\) in
\(L^{4/3}(Q)\).
\end{proposition}

\begin{proof}
For every smooth spatial test \(\varphi\), the continuity equation and
Cauchy--Schwarz give
\begin{equation}\label{eq:narrow-mod}
 |\langle\varphi,m_\nu(t)-m_\nu(s)\rangle|
 \le \nu|t-s|\|\Delta\varphi\|_\infty
 +\|D\varphi\|_\infty\sqrt{|t-s|}
              \left(\int_Q\frac{|j_\nu|^2}{m_\nu}\right)^{1/2}.
\end{equation}
The space of probability measures on the compact torus is compact for narrow
convergence. Diagonal extraction against a countable dense family of
spatial tests therefore yields a continuous limiting probability curve,
with convergence uniform in time. Extract also weak limits in \(L^2\)
for the densities and \(L^{4/3}\) for the fluxes. One may obtain these
by diagonalizing bounded pairings against a countable dense family of
dual tests, then using \(L^p\) duality
\cite[Theorem~6.15, p.~190]{Folland}. Time integration identifies the weak
density with the narrow curve at almost every time.
Passing to \eqref{eq:weak-ce} proves
\((m,j)\in\mathcal K_0\) once lower semicontinuity gives finite action.
The diffusion term tends to zero since \(\nu m_\nu\to0\) in \(L^2\);
the initial marginal remains exactly \(m_0\).

Both nonnegative action terms are weakly lower semicontinuous. For the
kinetic term this follows from
\[
 \frac{|k|^2}{2\rho}
       =\sup_b\left(k\cdot b-\frac{\rho|b|^2}{2}\right).
\]
The integral equals the supremum over smooth bounded fields. To check
this without a positivity assumption, take \(b_N\) in the direction
of \(k\), of magnitude \(\min(|k|/\rho,N)\) where \(\rho>0\),
and magnitude \(N\) where \(\rho=0,k\ne0\).
Its nonnegative dual integrand increases to the perspective, including
the infinite vacuum case. Smooth convolution approximates \(b_N\)
in \(L^1\) and along a subsequence almost everywhere, keeping the
bound \(N\). Dominated convergence applies with
\(N|k|+N^2\rho/2\), and monotone convergence lets \(N\to\infty\).
Each smooth dual pairing is weakly continuous, proving the assertion.
The same argument with \(c_N=\min(\rho,N)\) and
\(\rho^2/2=\sup_c(\rho c-c^2/2)\) applies to the density square.
The convolution approximation is the usual approximate-identity
property \cite[Theorem~8.14, pp.~242--243]{Folland}; periodic functions are localized
by a compact cutoff, and time extensions are zero. Thus
\begin{equation}\label{eq:liminf}
 J_0(m,j)\le\liminf_{\nu\downarrow0}I_\nu
\end{equation}
along the chosen subsequence, with the terminal term converging narrowly.

For the opposite inequality, take any finite-action
\((\rho,k)\in\mathcal K_0\) and set
\begin{equation}\label{eq:recovery}
 \rho_\nu(t)=P_{\nu t}\rho(t),\qquad
 k_\nu(t)=P_{\nu t}k(t).
\end{equation}
Use \(P_{\nu t}\psi(t)\) as the test in the zero-diffusion equation.
Since
\(\partial_t(P_{\nu t}\psi)=P_{\nu t}(\psi_t+\nu\Delta\psi)\),
self-adjointness and commutation with \(D\) give
\eqref{eq:weak-ce} for \((\rho_\nu,k_\nu)\) with diffusivity \(\nu\).
At time zero \(P_0=\mathrm{id}\), so \(\rho_\nu(0)=m_0\) exactly.
For the positive probability kernel, weighted Cauchy--Schwarz gives
\[
 |P_s k|^2\le P_s\rho\,P_s(|k|^2/\rho),\qquad
 (P_s\rho)^2\le P_s(\rho^2).
\]
Integration proves that neither integral action term increases.
Moreover
\(\int g\,d\rho_\nu(T)=\int P_{\nu T}g\,d\rho(T)\to\int g\,d\rho(T)\).
Consequently
\[
 \limsup_{\nu\downarrow0}I_\nu\le I_0 .
\]
The set \(\mathcal K_0\) is nonempty, since \((m_0,0)\) is admissible.
Together with \eqref{eq:liminf}, this proves attainment and convergence
of the minima.

If two minimizers had different densities, the density-square term
at their midpoint would decrease the average action by
\(\frac18\int(\rho-\widetilde\rho)^2\), while the kinetic perspective
is convex. Thus the densities coincide. At a common density the
kinetic midpoint defect is
\(\frac18\int|k-\widetilde k|^2/\rho\), and both fluxes are zero
in vacuum. This proves uniqueness. Equality of the space-time densities
also identifies their narrow curves at all times by continuity.
Finally, convergence of total action and separate lower semicontinuity
force convergence of the density-square integrals. Weak convergence
and this norm convergence imply strong \(L^2\) convergence. Every
vanishing sequence has the same conclusion, so the density convergence
and convergence of the minima hold for the entire family.
\end{proof}

Let \(u\) be any uniform accumulation point on the closed cylinder supplied
by Proposition~\ref{prop:modulus}. Passing in \eqref{eq:dual-value} gives
\begin{equation}\label{eq:limit-dual}
 \int_{\T^d}m_0u(0)=I_0+\frac12\int_Qm^2 .
\end{equation}
The gradient estimate \eqref{eq:Du} identifies \(Du\in L^2(Q)\) as the
weak limit of a further subsequence of gradients. Convex lower
semicontinuity against nonnegative tests in the Hamilton--Jacobi equation
then gives
\begin{equation}\label{eq:subdist}
 -u_t+\tfrac12|Du|^2\le m \quad\hbox{in distributions on }Q .
\end{equation}
The term \(\nu\Delta u_\nu\) tends to zero distributionally, and
the terminal trace is \(u(T)=g\). Both \eqref{eq:limit-dual} and
\eqref{eq:subdist} hold for every value accumulation point with the same
pair \((m,j)\).

\subsection{A common family of calibrated paths}

We use a path representation of the finite-action continuity equation.
Its construction and the subsequent calibration are the quadratic
specializations of the path arguments in \cite[Section~4.3]{Cardaliaguet}.
We give the details needed here, without invoking that paper's general
MFG existence or uniqueness theorem.

\begin{lemma}\label{lem:paths}
There is a probability law \(\eta\) on
\(\Gamma=C([0,T];\T^d)\), with the uniform metric, such that
\begin{equation}\label{eq:path-law}
 (e_t)_\#\eta=m(t)\quad(0\le t\le T),\qquad
 \int_\Gamma\int_0^T|\dot\gamma|^2\,ds\,d\eta(\gamma)
           \le\int_Q\frac{|j|^2}{m}.
\end{equation}
Here \(e_t(\gamma)=\gamma(t)\). The law is concentrated on absolutely
continuous paths with square-integrable velocities, using local
Euclidean lifts on the torus.
\end{lemma}

\begin{proof}
For \(\delta>0\), smooth only in space with the torus heat semigroup
\(P_\delta=e^{\delta\Delta}\):
\(m^\delta=P_\delta m\), \(j^\delta=P_\delta j\).
The positive heat kernel and mass one give
\(m^\delta(t,x)\ge\kappa_\delta>0\).
All spatial derivatives of \(m^\delta\) are uniformly bounded in time,
and those of \(j^\delta(t)\) are bounded by
\(C_\delta\|j(t)\|_1\). Thus
\(b^\delta=j^\delta/m^\delta\) has integrable-in-time supremum
and spatial Lipschitz bounds.
Its integral equation has a unique flow, by Picard iteration on finitely
many intervals where the integrated Lipschitz bound is below \(1/2\).
Reversing time supplies the inverse flow.

The flow transports \(m^\delta(0)\) to \(m^\delta(t)\). Indeed,
\(\partial_tm^\delta=-\operatorname{div}j^\delta\) has an
integrable-in-time smooth-spatial right side. A smooth terminal test
composed with the backward flow satisfies
\(\zeta_t+b^\delta\cdot D\zeta=0\); differentiating
\(\int\zeta m^\delta\) gives zero. The same identity along the flow
identifies every terminal test integral of the push-forward.
The required time derivatives are absolutely continuous, so the product
rule follows by integration of their \(L^1\) derivatives.

Let \(\eta^\delta\) be the law of these flow paths, with initial
distribution \(m^\delta(0)\). Weighted Cauchy--Schwarz for the heat
kernel yields
\[
 \int_\Gamma\int_0^T|\dot\gamma|^2ds\,d\eta^\delta
 =\int_Q\frac{|j^\delta|^2}{m^\delta}
 \le\int_Q\frac{|j|^2}{m}.
\]
The path energy, infinite outside the finite-energy paths, is lower
semicontinuous and has compact sublevels. One convenient expression is
\[
 \mathcal E(\gamma)=
 \sup_{0=t_0<\cdots<t_n=T}
 \sum_i\frac{\operatorname{dist}(\gamma(t_i),\gamma(t_{i-1}))^2}
                    {t_i-t_{i-1}} .
\]
For an absolutely continuous path, Cauchy--Schwarz bounds this expression
by its energy. Conversely, polygonal interpolants on vanishing-mesh
partitions have bounded \(L^2\) velocity whenever the displayed supremum
is finite. On local torus charts, weak extraction of these velocities
and uniform convergence of the interpolants give an absolutely continuous
limit lift. Interval averages of its velocity converge in \(L^2\),
so the supremum equals its energy. This also proves lower
semicontinuity, since each partition expression is continuous.
Its sublevels have modulus
\(\operatorname{dist}(\gamma(t),\gamma(s))\le\sqrt{R|t-s|}\)
at energy \(R\), and are compact by Arzel\`a--Ascoli.

Markov's inequality therefore makes the laws tight.
The path space is complete and separable, so Prokhorov compactness
and the lower-semicontinuity inequality for integrals give a narrow
limit \(\eta\) with the asserted energy bound; see
\cite[Appendix~C, Theorems~C.9 and C.12, p.~430 of the author manuscript]{AGZ}.
The evaluation maps are continuous and \(P_\delta m(t)\to m(t)\)
narrowly for every time. They give all the marginals in
\eqref{eq:path-law}.
\end{proof}

\begin{proposition}\label{prop:calibration}
Every value accumulation point \(u\) satisfies, on one \(\eta\)-full set,
\begin{equation}\label{eq:calibrated}
 u(t,\gamma(t))=g(\gamma(T))+
 \int_t^T\left(\tfrac12|\dot\gamma(s)|^2+
                          m(s,\gamma(s))\right)ds,\qquad 0\le t\le T.
\end{equation}
In particular, any two value accumulation points \(u,v\) agree
\(m(t)\)-almost everywhere at every time, and agree almost everywhere
on the occupied region \(\{m>0\}\subset Q\).
\end{proposition}

\begin{proof}
Convolve \eqref{eq:subdist} with a positive smooth unit-mass
space-time kernel. Denote the resulting functions by \(u_\epsilon\),
\(m_\epsilon\); these differ from the spatial heat regularizations
used in Lemma~\ref{lem:paths}. Convexity gives
\[
 -u_{\epsilon,t}+\tfrac12|Du_\epsilon|^2\le m_\epsilon
\]
on each compact interior time interval for small \(\epsilon\).
Along an absolutely continuous finite-energy path the chain rule and
completion of a square give
\[
 -\frac d{ds}u_\epsilon(s,\gamma(s))
 \le m_\epsilon-\tfrac12|Du_\epsilon|^2
                      -Du_\epsilon\cdot\dot\gamma
 \le m_\epsilon+\tfrac12|\dot\gamma|^2 .
\]
The endpoints pass by local uniform convergence of \(u_\epsilon\).
The running cost passes for almost every path because
\begin{align}
 \int_\Gamma\int_0^T
 |m_\epsilon(s,\gamma(s))-m(s,\gamma(s))|\,ds\,d\eta
 &=\int_Qm|m_\epsilon-m| \notag\\
 &\le\|m\|_2\|m_\epsilon-m\|_2\longrightarrow0 .
 \label{eq:family-trace}
\end{align}
Here the time extension of \(m\) is zero; periodic spatial convolution
is justified by localization of the extension. Restrict the smooth
inequality to interior intervals before taking the limit.
A subsequence gives convergence of the absolute path integrals almost
surely. Pass first for rational endpoints, then for all endpoints by
continuity and integrability of the path cost, finally including
zero and \(T\) by the value traces. This gives, on one full set,
\begin{equation}\label{eq:path-ineq}
 u(a,\gamma(a))\le u(b,\gamma(b))+
 \int_a^b\left(\tfrac12|\dot\gamma|^2+m(s,\gamma(s))\right)ds
 \quad(0\le a<b\le T).
\end{equation}
The cost is integrable: its expectation is at most
\(\frac12\int |j|^2/m+\int m^2\).
A Borel representative of \(m\) may be chosen arbitrarily; the
occupation identity in \eqref{eq:family-trace} makes changes on
space-time null sets irrelevant to almost every path-time.

For \(a=0,b=T\), the expected right side of \eqref{eq:path-ineq}
is at most
\[
 \frac12\int_Q\frac{|j|^2}{m}+\int_Qm^2+\int g\,dm(T)
 =I_0+\frac12\int_Qm^2
 =\int_{\T^d}m_0u(0),
\]
by \eqref{eq:limit-dual}. The inequality supplies the reverse bound,
so its nonnegative pathwise gap has zero mean and vanishes almost surely.
Split that zero gap at \(t\). Both partial gaps are nonnegative and
therefore vanish, proving \eqref{eq:calibrated} for every \(t\)
on the same full set.

Every other accumulation point satisfies the identical formula for
the same \(m,j,\eta\). Intersect the two full sets and push forward
by \(e_t\) to get equality \(m(t)\)-almost everywhere.
For almost every time this marginal has density \(m(t,x)\);
Fubini then gives \(\int_Q m|u-v|=0\), proving the last assertion.
\end{proof}

\section{Uniqueness of the value limit}\label{sec:vacuum}

The equality on the occupied region leaves only vacuum to consider.
Every accumulation point is a viscosity supersolution of the homogeneous
equation
\begin{equation}\label{eq:super}
 -u_t+\tfrac12|Du|^2\ge0 .
\end{equation}
For completeness, if a smooth test \(\phi\) touches \(u\) from below
at an interior point, make the touch strict and take nearby interior
minimizers of \(u_\nu-\phi\). There
\(u_{\nu,t}=\phi_t\), \(Du_\nu=D\phi\) and
\(\Delta u_\nu\ge\Delta\phi\). Hence
\[
 -\phi_t-\nu\Delta\phi+\tfrac12|D\phi|^2\ge m_\nu\ge0 .
\]
Uniform convergence permits passage to the touching point, proving
\eqref{eq:super}. A fourth-power local perturbation removes the
strict-touch restriction. No pointwise limiting density is used.

Let \(u,v\) be two accumulation points and set
\(O=\{(t,x)\in Q:u(t,x)>v(t,x)\}\).
Proposition~\ref{prop:calibration} gives \(m=0\) almost everywhere in \(O\).
Equation \eqref{eq:subdist} therefore makes \(u\) a distributional
subsolution of the homogeneous Hamilton--Jacobi equation there.
It is also a viscosity subsolution. Indeed, on compact subsets of \(O\),
positive convolution and convexity give smooth subsolutions converging
uniformly to \(u\). Interior maxima against a smooth upper test pass
the subsolution inequality to the limit. The gradient \(Du\in L^2\)
justifies this convolution.

We show directly that \(O\) is empty. This is the homogeneous comparison
used for the vacuum extension in \cite[proof of Theorem~4.2]{Cardaliaguet}.
Suppose \(u(t_0,x_0)>v(t_0,x_0)\). Choose \(0<a<t_0\) and small
\(\varepsilon>0\) so that
\[
 w(t,x)=u(t,x)-\frac{\varepsilon}{t-a}
\]
still exceeds \(v\) at \((t_0,x_0)\).
On \(O\cap\{t>a\}\) it is a strict viscosity subsolution:
\[
 -w_t+\tfrac12|Dw|^2\le-\frac{\varepsilon}{(t-a)^2}.
\]
The positive maximum of \(w-v\) stays away from \(a\), where \(w\)
tends below \(v\), from \(T\), where \(u=v=g\), and from the
interior boundary of \(O\), where \(u=v\).
On a compact time interval containing this maximum, maximize
\[
 w(t,x)-v(s,y)
 -\frac{\operatorname{dist}(x,y)^2+|t-s|^2}{2\delta}.
\]
Boundedness and uniform continuity force the separation to zero as
\(\delta\downarrow0\), while the limiting difference remains positive.
The time coordinates are therefore interior, and the \(w\)-test point
is in \(O\). For small \(\delta\), the shortest lift is unique and the
penalty is smooth there. The upper test for \(w\) and lower test for
\(v\) have the same derivatives
\[
 q=(t-s)/\delta,\qquad p=(x-y)/\delta
\]
in these lifts. They would give simultaneously
\[
 -q+\tfrac12|p|^2\le-\frac{\varepsilon}{(t-a)^2}<0,
 \qquad -q+\tfrac12|p|^2\ge0 ,
\]
a contradiction. The gradients need not be bounded: the Hamiltonian
is independent of time and position and cancels exactly.
No initial value condition for the first-order value function is
required. Thus \(u\le v\); interchanging them proves equality.

\begin{proof}[Proof of Theorem~\ref{thm:main}]
Every vanishing sequence has a uniformly convergent subsequence on
compact interior slabs, and its limit is continuous by
Proposition~\ref{prop:modulus}. The argument above identifies every
accumulation point with the same function \(u\). If convergence failed
on a fixed slab, a sequence with uniform error bounded below by a
positive number would admit a subsequence converging to that same \(u\),
a contradiction. The single function obtained this way works for every
\(0<\tau<T/2\), which proves the entire-family assertion.
\end{proof}

\section*{AI use}
AI tools were used in developing the argument, drafting and revising the
exposition, checking references and validating the document. Two additional
AI contexts reviewed a frozen draft for reader comprehension and mathematical
objections. The manuscript is submitted for independent mathematical review.

\begin{thebibliography}{99}

\bibitem{AGZ}
G.~W.~Anderson, A.~Guionnet and O.~Zeitouni,
\emph{An Introduction to Random Matrices},
Cambridge Studies in Advanced Mathematics~118,
Cambridge University Press, 2009.
\href{https://doi.org/10.1017/CBO9780511801334}{doi:10.1017/CBO9780511801334}.
An author manuscript containing Appendix~C is available at
\url{https://staff.math.su.se/shapiro/UIUC/RandomMatrices.pdf};
the Appendix~C locators refer to this manuscript.

\bibitem{Cardaliaguet}
P.~Cardaliaguet,
\emph{Weak solutions for first order mean field games with local coupling},
arXiv:1305.7015v1, 2013.
\url{https://arxiv.org/abs/1305.7015v1}.

\bibitem{CLLP}
P.~Cardaliaguet, J.-M.~Lasry, P.-L.~Lions and A.~Porretta,
Long time average of mean field games,
\emph{Networks and Heterogeneous Media} \textbf{7} (2012), no.~2, 279--301.
\href{https://doi.org/10.3934/nhm.2012.7.279}{doi:10.3934/nhm.2012.7.279}.

\bibitem{Folland}
G.~B.~Folland,
\emph{Real Analysis: Modern Techniques and Their Applications},
second edition, John Wiley \& Sons, Inc., 1999.

\bibitem{GraberMeszaros}
P.~Jameson Graber and Alp\'ar R.~M\'esz\'aros,
Sobolev regularity for first order mean field games,
\emph{Annales de l'Institut Henri Poincar\'e C, Analyse non lin\'eaire}
\textbf{35} (2018), no.~6, 1557--1576.
\href{https://doi.org/10.1016/j.anihpc.2018.01.002}{doi:10.1016/j.anihpc.2018.01.002}.

\bibitem{Lawler}
G.~F.~Lawler,
\emph{Stochastic Calculus: An Introduction with Applications},
author manuscript, revised 15 February 2023.
\url{https://www.math.uchicago.edu/~lawler/finbook.pdf}.

\bibitem{Nirenberg}
L.~Nirenberg,
On elliptic partial differential equations,
\emph{Annali della Scuola Normale Superiore di Pisa, Classe di Scienze},
series~3, \textbf{13} (1959), no.~2, 115--162.
\url{https://www.numdam.org/item/ASNSP_1959_3_13_2_115_0/}.

\bibitem{TangZhang}
W.~Tang and Y.~P.~Zhang,
The convergence rate of vanishing viscosity approximations for mean field games,
\emph{SIAM Journal on Mathematical Analysis} \textbf{57} (2025),
no.~3, 3217--3254.
\href{https://doi.org/10.1137/24M1640008}{doi:10.1137/24M1640008}.

\end{thebibliography}
\end{document}
